\subsection{INVERSE SYSTEMS OF MAPPINGS}

PROPOSITION 1. Let I be an ordered set, let \((\mathbf{E}_{\alpha}, f_{\alpha \beta})\) be an inverse system of sets relative to I, let \(\mathbf{E} = \lim \mathbf{E}_{\alpha}\) be its inverse limit, and for each \(\alpha \in I\) let

\[
f _ {\alpha}: \mathbf {E} \rightarrow \mathbf {E} _ {\alpha}
\]

be the canonical mapping. For each \(\alpha \in \mathbf{I}\) , let \(u_{\alpha}\) be a mapping of a set \(\mathbf{F}\) into \(\mathbf{E}_{\alpha}\) such that

\[
f _ {\alpha \beta} \circ u _ {\beta} = u _ {x}\tag{5}
\]

whenever \(\alpha \leqslant \beta\) .

Then :

\begin{enumerate}
\def\labelenumi{(\alph{enumi})}
\tightlist
\item
  there exists a unique mapping \(u\) of \(\mathbf{F}\) into \(\mathbf{E}\) such that
\end{enumerate}

\[
u _ {\alpha} = f _ {\alpha} \circ u\tag{6}
\]

\[
\text { for   all } \alpha \in \mathbf {I};
\]

\begin{enumerate}
\def\labelenumi{(\alph{enumi})}
\setcounter{enumi}{1}
\tightlist
\item
  the mapping \(u\) is injective if and only if, for each pair of distinct elements \(y, z\) of \(\mathbf{F}\) , there exists \(\alpha \in \mathbf{I}\) such that \(u_{\alpha}(y) \neq u_{\alpha}(z)\) .
\end{enumerate}

For the relation \(u_{\alpha} = f_{\alpha} \circ u\) means that for each \(y \in F\) we have

\[
\operatorname{pr} _ {\alpha} (u (\gamma)) = u _ {\alpha} (\gamma);
\]

the element \(u(y) \in \prod_{\alpha \in I} E_{\alpha}\) is uniquely determined by \(u(y) = (u_{\alpha}(y))_{\alpha \in I}\) . It remains to be shown that \(u(y) \in E\) for all \(y \in F\) , in other words, that

\[
\operatorname{pr} _ {\alpha} (u (y)) = f _ {\alpha \beta} (\operatorname{pr} _ {\beta} (u (y)))
\]

whenever \(\alpha \leqslant \beta\) . But this can be written in the form

\[
u _ {\alpha} (y) = f _ {\alpha \beta} (u _ {\beta} (y))
\]

and therefore follows from (5). The second part of the Proposition follows immediately from the definitions.

COROLLARY 1. Let \((\mathbf{E}_{\alpha}, f_{\alpha \beta})\) and \((\mathbf{F}_{\alpha}, g_{\alpha \beta})\) be two inverse systems of sets relative to the same index set \(\mathbf{I}\) ; let \(\mathbf{E} = \lim \mathbf{E}_{\alpha}\) , \(\mathbf{F} = \lim \mathbf{F}_{\alpha}\) , and let \(f_{\alpha}\) (resp. \(g_{\alpha}\) ) be the canonical mapping of \(\mathbf{E}\) into \(\overleftarrow{\mathbf{E}_{\alpha}}\) (resp. of \(\overleftarrow{\mathbf{F}}\) into \(\mathbf{F}_{\alpha}\) ) for each \(\alpha \in \mathbf{I}\) , For each \(\alpha \in \mathbf{I}\) , let \(u_{\alpha}\) be a mapping of \(\mathbf{E}\) into \(\mathbf{F}_{\alpha}\) such that the diagram

\[
\begin{array}{c} \mathbf {E} _ {\beta} \xrightarrow {u _ {\beta}} \mathbf {F} _ {\beta} \\ f _ {\alpha \beta} \Big \downarrow \qquad \qquad \qquad \qquad \Big \downarrow g _ {\alpha \beta} \\ \mathbf {E} _ {\alpha} \xrightarrow [ u _ {\alpha} ]{} \mathbf {F} _ {\alpha} \end{array}
\]

is commutative (*) whenever \(\alpha \leqslant \beta\) . Then there exists a unique mapping \(u: \mathbf{E} \to \mathbf{F}\) such that for each \(\alpha \in \mathbf{I}\) the diagram

\[
\begin{array}{c} \mathbf {E} \xrightarrow {u} \mathbf {F} \\ f _ {\alpha} \Big \downarrow \qquad \qquad \qquad \qquad \qquad \qquad \qquad \qquad \qquad \Big \downarrow g _ {\alpha} \\ \mathbf {E} _ {\alpha} \xrightarrow [ u _ {\alpha} ]{} \mathbf {F} _ {\alpha} \end{array}
\]

is commutative.

Put \(v_{\alpha} = u_{\alpha} \circ f_{\alpha}\) . If \(\alpha \leqslant \beta\) , we have, by (2),

\[
g _ {\alpha \beta} \circ v _ {\beta} = g _ {\alpha \beta} \circ u _ {\beta} \circ f _ {\beta} = u _ {\alpha} \circ f _ {\alpha \beta} \circ f _ {\beta} = u _ {\alpha} \circ f _ {\alpha} = v _ {\alpha},
\]

and we may therefore apply Proposition 1 to the mappings \(v_{\alpha}\) ; hence the

existence and uniqueness of a mapping \(u: E \rightarrow F\) such that

\[
g _ {\alpha} \circ u = v _ {\alpha} = u _ {\alpha} \circ f _ {\alpha}
\]

for each \(\alpha \in I\) .

A family of mappings \(u_{\alpha}: \mathbf{E}_{\alpha} \to \mathbf{F}_{\alpha}\) which satisfies the conditions of Corollary 1 is called an inverse system of mappings of \((\mathbf{E}_{\alpha}, f_{\alpha\beta})\) into \((\mathbf{F}_{\alpha}, g_{\alpha\beta})\) . The mapping \(u\) defined in Corollary 1 is called the inverse limit of the family \((u_{\alpha})\) and is written \(u = \lim_{\leftarrow} u_{\alpha}\) when there is no risk of confusion.

COROLLARY 2. Let \((\mathbf{E}_{\alpha}, f_{\alpha \beta})\) , \((\mathbf{F}_{\alpha}, g_{\alpha \beta})\) , \((\mathbf{G}_{\alpha}, h_{\alpha \beta})\) be three inverse systems of sets relative to the same index set I; let \(\mathbf{E} = \varprojlim \mathbf{E}_{\alpha}\) , \(\mathbf{F} = \varprojlim \mathbf{F}_{\alpha}\) , \(\mathbf{G} = \varprojlim \mathbf{G}_{\alpha}\) , and let \(f_{\alpha}\) (resp. \(g_{\alpha}, h_{\alpha}\) ) be the canonical mapping of \(\mathbf{E}\) (resp. \(\mathbf{F}, \mathbf{G}\) ) into \(\mathbf{E}_{\alpha}\) (resp. \(\mathbf{F}_{\alpha}, \mathbf{G}_{\alpha}\) ). If \((u_{\alpha})\) and \((v_{\alpha})\) are two inverse systems of mappings, \(u_{\alpha}: \mathbf{E}_{\alpha} \to \mathbf{F}_{\alpha}\) , \(v_{\alpha}: \mathbf{F}_{\alpha} \to \mathbf{G}_{\alpha}\) , then the composite mappings \(v_{\alpha} \circ u_{\alpha}: \mathbf{E}_{\alpha} \to \mathbf{G}_{\alpha}\) form an inverse system of mappings, and we have

\[
\varprojlim (v _ {\alpha} \circ u _ {\alpha}) = (\varprojlim v _ {\alpha}) \circ (\varprojlim u _ {\alpha}).\tag{7}
\]

For if we put \(w_{\alpha} = v_{\alpha} \circ u_{\alpha}\) , then if \(\alpha \leqslant \beta\) we have

\[
w _ {\alpha} \circ f _ {\alpha \beta} = v _ {\alpha} \circ (u _ {\alpha} \circ f _ {\alpha \beta}) = v _ {\alpha} \circ (g _ {\alpha \beta} \circ u _ {\beta}) = (h _ {\alpha \beta} \circ v _ {\beta}) \circ u _ {\beta} = h _ {\alpha \beta} \circ w _ {\beta},
\]

which shows that \((w_{\alpha})\) is an inverse system of mappings. Furthermore, if \(u = \lim u_{\alpha}\) and \(v = \lim v_{\alpha}\) , then \(h_{\alpha} \circ (v \circ u) = (v_{\alpha} \circ g_{\alpha}) \circ u = (v_{\alpha} \circ u_{\alpha}) \circ f_{\alpha}\) for each \(\alpha \in I\) , and therefore, by the uniqueness of the inverse limit, we have \(v \circ u = \lim w_{\alpha}\) .

Let \((\mathbf{E}_{\alpha}, f_{\alpha \beta})\) be an inverse system of sets, and for each \(\alpha \in I\) let \(\mathbf{M}_{\alpha}\) be a subset of \(\mathbf{E}_{\alpha}\) . If \(f_{\alpha \beta}(\mathbf{M}_{\beta}) \subset \mathbf{M}_{\alpha}\) whenever \(\alpha \leqslant \beta\) , the \(\mathbf{M}_{\alpha}\) are said to form an inverse system of subsets of the \(\mathbf{E}_{\alpha}\) . Let \(g_{\alpha \beta}\) be the mapping of \(\mathbf{M}_{\beta}\) into \(\mathbf{M}_{\alpha}\) (where \(\alpha \leqslant \beta\) ) whose graph is the same as that of the restriction of \(f_{\alpha \beta}\) to \(\mathbf{M}_{\beta}\) . Then it is clear that \((\mathbf{M}_{\alpha}, g_{\alpha \beta})\) is an inverse system of sets and that

\[
\varprojlim \mathbf {M} _ {\alpha} = (\varprojlim \mathbf {E} _ {\alpha}) \cap \prod_ {\alpha \in \mathbf {I}} \mathbf {M} _ {\alpha}.\tag{8}
\]

PROPOSITION 2. Let \((\mathbf{E}_{\alpha}, f_{\alpha \beta})\) and \((\mathbf{E}_{\alpha}^{\prime}, f_{\alpha \beta}^{\prime})\) be two inverse systems of sets relative to I, and let \(u_{\alpha}\) be a mapping of \(\mathbf{E}_{\alpha}\) into \(\mathbf{E}_{\alpha}^{\prime}\) for each \(\alpha \in \mathbf{I}\) , such that the \(u_{\alpha}\) form an inverse system of mappings. Let \(u = \varprojlim u_{\alpha}\) . Then for each \(x^{\prime} = (x_{\alpha}^{\prime}) \in \mathbf{E}^{\prime} = \varprojlim \mathbf{E}_{\alpha}^{\prime}\) , the \(-\overleftarrow{u}_{\alpha}(x_{\alpha}^{\prime})\) form an inverse system of subsets of the \(\mathbf{E}_{\alpha}\) , and \(-\overleftarrow{u}(x^{\prime}) = \varprojlim -\overleftarrow{u}_{\alpha}(x_{\alpha}^{\prime})\) .

For if \(\alpha \leqslant \beta\) and \(x_{\beta} \in \overline{u}_{\beta}^{1}(x_{\beta}')\) , we have

\[
u _ {\alpha} (f _ {\alpha \beta} (x _ {\beta})) = f _ {\alpha \beta} ^ {\prime} (u _ {\beta} (x _ {\beta})) = f _ {\alpha \beta} ^ {\prime} (x _ {\beta} ^ {\prime}) = x _ {\alpha} ^ {\prime},
\]

from which the first assertion follows; and to say that \(x = (x_{\alpha}) \in \mathbf{E} = \lim \mathbf{E}_{\alpha}\) is such that \(u(x) = x'\) means, by definition, that \(u_{\alpha}(x_{\alpha}) = x_{\alpha}'\) for each \(\alpha \in I\) .

COROLLARY. If \(u_{\alpha}\) is injective (resp. bijective) for each \(\alpha \in I\) , then \(u\) is injective (resp. bijective).

With the notation of Proposition 2, the images \(u_{\alpha}(\mathbf{E}_{\alpha})\) also form an inverse system of subsets of the \(\mathbf{E}_{\alpha}'\) , and we have

\[
u (\mathrm{E}) \subset \varprojlim u _ {\alpha} (\mathrm{E} _ {\alpha});\tag{9}
\]

but the two sides of this relation are not necessarily equal (Exercise 4).

PROPOSITION 3. Let I be a preordered set, let \((\mathbf{E}_{\alpha}, f_{\alpha \beta})\) be an inverse system of sets relative to I, and let \(\mathbf{E} = \lim \mathbf{E}_{\alpha}\) . Let J be a cofinal subset of I such that J is right directed, and let \(\mathbf{E}'\) be the inverse limit of the inverse system of sets obtained from \((\mathbf{E}_{\alpha}, f_{\alpha \beta})\) by restricting the index set to J. Then the canonical mapping \(g\) of E into \(\mathbf{E}'\) (no. 1, formula (3)) is bijective.

For each \(\alpha \in J\) , let \(f_{\alpha}'\) be the canonical mapping \(\mathbf{E}' \to \mathbf{E}_{\alpha}\) . Then, by (2) and (5), \(g\) is the unique mapping of \(\mathbf{E}\) into \(\mathbf{E}'\) such that \(f_{\alpha} = f_{\alpha}' \circ g\) for all \(\alpha \in J\) (Proposition 1). We shall show that \(g\) is injective by using the criterion of Proposition 1. If \(x, y\) are distinct elements of \(\mathbf{E}\) , then by definition there exists \(\alpha \in I\) such that \(f_{\alpha}(x) \neq f_{\alpha}(y)\) ; since \(J\) is cofinal in \(I\) , there exists \(\lambda \in J\) such that \(\alpha \leqslant \lambda\) ; since \(f_{\alpha\lambda}(f_{\lambda}(x)) \neq f_{\alpha\lambda}(f_{\lambda}(y))\) , we have \(f_{\lambda}(x) \neq f_{\lambda}(y)\) . It remains to be shown that \(g\) is surjective. Let \(x' = (x_{\lambda}')_{\lambda \in J}\) be an element of \(\mathbf{E}'\) . For each \(\alpha \in I\) there exists \(\lambda \in J\) such that \(\alpha \leqslant \lambda\) , and the element \(f_{\alpha\lambda}(x_{\lambda}')\) does not depend on the index \(\lambda \in J\) such that \(\alpha \leqslant \lambda\) ; for if \(\mu \in J\) is such that \(\alpha \leqslant \mu\) , there exists \(\nu \in J\) such that \(\lambda \leqslant \nu\) and \(\mu \leqslant \nu\) , hence \(f_{\alpha\lambda}(x_{\lambda}') = f_{\alpha\lambda}(f_{\lambda\nu}(x_{\nu}') = f_{\alpha\nu}(x_{\nu}')\) , and similarly \(f_{\alpha\mu}(x_{\mu}') = f_{\alpha\nu}(x_{\nu}')\) . Let \(x_{\alpha}\) be the common value of the \(f_{\alpha\lambda}(x_{\lambda}')\) for the \(\lambda \in J\) such that \(\alpha \leqslant \lambda\) , and let \(x = (x_{\alpha})_{\alpha \in I}\) . Then \(x \in E\) , because if \(\alpha \leqslant \beta\) and if \(\lambda \in J\) is such that \(\beta \leqslant \lambda\) , we have

\[
f _ {\alpha \beta} (x _ {\beta}) = f _ {\alpha \beta} (f _ {\beta \lambda} (x _ {\lambda} ^ {\prime})) = f _ {\alpha \lambda} (x _ {\lambda} ^ {\prime}) = x _ {\alpha}.
\]

Finally, we have \(x_{\lambda}' = f_{\lambda \lambda}(x_{\lambda}')\) for all \(\lambda \in \mathbf{J}\) , hence \(x_{\lambda} = x_{\lambda}'\) for all \(\lambda \in \mathbf{J}\) ; in other words, \(f(x_{\lambda}) = x_{\lambda}',\) so that \(g(x) = x'\) . Therefore \(g\) is surjective and the proof is complete.

In particular, if I has a greatest element \(\omega\) , we may take \(J = \{\omega\}\) , so that \(\lim E_{\alpha}\) is canonically identified with \(E_{\omega}\) .

Remarks

\begin{enumerate}
\def\labelenumi{(\arabic{enumi})}
\tightlist
\item
  For each \(\alpha \in I\) put \(E_{\alpha}' = f_{\alpha}(E)\) . Then the sets \(E_{\alpha}'\) form an inverse system of subsets of the \(E_{\alpha}\) by reason of (2), and it is immediately clear that \(\lim E_{\alpha}' = E = \lim E_{\alpha}\) . The mapping \(f_{\alpha\beta}' : E_{\beta}' \to E_{\alpha}'\) (where \(\alpha \leqslant \beta\) ), whose graph is the same as that of the restriction of \(f_{\alpha\beta}\) to \(E_{\beta}\) , is surjective, and we have
\end{enumerate}

\[
\mathbf {E} _ {\alpha} ^ {\prime} = f _ {\alpha} (\mathbf {E}) \subset \bigcap_ {\beta \geqslant \alpha} f _ {\alpha \beta} (\mathbf {E} _ {\beta})\tag{10}
\]

for all \(\alpha \in I\) .

\begin{enumerate}
\def\labelenumi{(\arabic{enumi})}
\setcounter{enumi}{1}
\item
  Let I be a (right) directed ordered set, let \((\mathbf{E}_{\alpha}, f_{\alpha\beta})\) be an inverse system of sets relative to I, and for each \(\alpha \in I\) let \(u_{\alpha}: F \to E_{\alpha}\) be a mapping such that the family \((u_{\alpha})\) satisfies formula (5). Consider the inverse system \((\mathbf{F}_{\alpha}, i_{\alpha\beta})\) indexed by I, where \(F_{\alpha} = F\) for all \(\alpha \in I\) and \(i_{\alpha\beta}\) is the identity mapping of F. Then (no. 1, Example 2) F is canonically identified with \(\lim F_{\alpha}\) . If we consider \(u_{\alpha}\) as a mapping of \(F_{\alpha}\) into \(E_{\alpha}\) , then \((u_{\alpha})\) is an inverse system of mappings, and the mapping \(u: F \to E\) defined by (6) is identified with the inverse limit of this system of mappings. Hence, by abuse of language, we write \(u = \lim u_{\alpha}\) .
\item
  Let I be an ordered set and let \((\mathbf{E}_{\alpha}, f_{\alpha \beta})\) be an inverse system of sets relative to I. For each finite subset \(J\) of \(I\) , let \(F_{J}\) be the inverse limit of the (finite) inverse system obtained from \((\mathbf{E}_{\alpha}, f_{\alpha \beta})\) by restricting the index set to J. If J and K are any two finite subsets of I such that \(J \subset K\) , let \(g_{JK}\) denote the canonical mapping (3) of \(F_{K}\) into \(F_{J}\) . Then the relation (4) shows that \((F_{J}, g_{JK})\) is an inverse system of sets relative to the directed set (with respect to the relation \(c\) ) \(\mathfrak{F}(I)\) of finite subsets of I. Next, for each \(J \in \mathfrak{F}(I)\) let \(h_{J}: E \to F_{J}\) be the canonical mapping (3). By virtue of (4) and with the abuse of language mentioned in Remark (2), \((h_{J})\) is an inverse system of mappings. Put \(h = \lim_{\leftarrow} h_{J}: E \to F = \lim_{\leftarrow} F_{J}\) , and let us show that \(h\) is a bijection (called canonical). Indeed, let \(y = (y_{J}) \in F\) . By definition we have \(y_{J} = (x_{\alpha, J})_{\alpha \in J}\) , where \(x_{\alpha, J} \in E_{\alpha}\) for all \(\alpha \in J\) . If \(J \subset K\) , then by definition of the mapping \(g_{JK}\) and because \(y_{J} = g_{JK}(y_{K})\) , we have \(x_{\alpha, J} = x_{\alpha, K}\) for all \(\alpha \in J\) . Hence, given \(\alpha \in I\) , there is a unique element \(x_{\alpha} \in E_{\alpha}\) such that \(x_{\alpha} = x_{\alpha, J}\) for all finite subsets J of I which contain \(\alpha\) . If \(\alpha \leqslant \beta\) , there is a finite subset J of I which contains both \(\alpha\) and \(\beta\) ; hence \(x_{\alpha} = f_{\alpha \beta}(x_{\beta})\) by definition. Consequently \(x = (x_{\alpha})\) is the unique element of E such that \(h(x) = y\) .
\end{enumerate}

